\begin{textsample}{POS Dim 7 – 1970 – Score 8.00 – 1970/tv\_com\_1970\_92.txt}  \label{ex:f7_pos_033}
A young adult woman with curly red hair stands against a dark \textbf{studio} background with a warm reddish glow behind her.
She wears a light button-up shirt with a narrow red tie and dark, shiny overall-style straps over the front.
Facing forward, she smiles slightly and speaks as if introducing her choice.
She then lifts an old-fashioned black \textbf{telephone} into view, holding it with both hands.
The phone has an antique look : a dark upright stand, a separate earpiece or \textbf{receiver}, and a round dial area on a \textbf{base}.
She turns her head toward it admiringly, then back toward the viewer, presenting it as if it reflects her personal taste.
The scene changes to an elderly woman with short, light-colored hair, also in a dark, warmly lit setting.
She laughs broadly while holding a fluffy cat close to her chest.
Beside her is a bright novelty \textbf{telephone} shaped around a Mickey Mouse figure : Mickey stands on a red \textbf{base} with yellow shoes, with a yellow handset set vertically along one side.
The cat looks forward while the woman cuddles it, and the playful phone sits prominently in front of them as one of her carefully chosen “ friends. ” Small white copyright text is visible near the bottom of the Mickey phone \textbf{base}.
A tabletop display follows, crowded with many \textbf{telephones} of different eras, colors, and shapes.
The arrangement is dense and theatrical, lit from above against a dark background.
At first, antique models dominate : wooden and brass-looking phones with rotary dials, upright candlestick \textbf{styles}, and polished black or dark brown handsets.
Mixed among them are more modern push-button desk phones in cream, white, tan, and brown.
A glossy red phone, a cream rotary phone, and the Mickey Mouse novelty phone are also visible among the collection.
The phones overlap across the table, suggesting a wide variety of personal \textbf{styles}, from decorative antiques to \textbf{sleek} contemporary designs.
As the display continues, the view settles more and more on a plain white push-button desk \textbf{telephone} near the center.
It has a rectangular sloped \textbf{base} with a grid of square buttons and a white handset resting across the top.
Around it are other phones : a tan phone to the right with its handset in place, a red phone behind it, a cream rotary phone farther back, and darker models to the left.
The white push-button phone becomes the visual focus, standing out as simpler and more standard than the more decorative phones surrounding it.
A hand reaches in at the top of the white phone and lifts away its \textbf{outer} shell or casing.
The smooth white exterior rises off, revealing that beneath the stylish surface is the \textbf{actual} working \textbf{telephone} mechanism : a compact keypad assembly with black square buttons, metal supports, exposed wiring, and internal components sitting in the \textbf{base}.
The contrast is clear visually—the attractive \textbf{outer} \textbf{form} is only the cover, while the essential functioning parts underneath are what make it a real \textbf{telephone}.
Finally, a bald adult man in a white suit jacket and white turtleneck appears against the same dark background with a warm amber halo behind him.
He smiles gently and holds up a distinctive white \textbf{telephone} shaped like a smooth circular ring or rounded sculptural \textbf{form}.
A white cord hangs down from one side.
He presents it at chest height with both hands, centered in front of him, as if displaying his own obvious choice.
The commercial ends with the \textbf{emphasis} on choosing a genuine Bell phone in a \textbf{style} that suits the individual.

% matched lemmas: actual, base, emphasis, form, outer, receiver, sleek, studio, style, telephone
\end{textsample}
